%  -----------------------------------------------------------------------------
%  Author         : Bimalka Piyaruwan Thalagala
%  GitHub         : https://github.com/bimalka98
%  Last Modified  : 2022-05-10
%  -----------------------------------------------------------------------------

\documentclass[a4paper,12pt]{report}%,twocolumn
%% packages

\usepackage{blindtext} % needed for creating dummy text passages
%\usepackage{ngerman} % needed for German default language
\usepackage{amsmath} % needed for command eqref
\usepackage{amssymb} % needed for math fonts
\usepackage[colorlinks=true,breaklinks]{hyperref} % needed for creating hyperlinks in the document, the option colorlinks=true gets rid of the awful boxes, breaklinks breaks lonkg links (list of figures), and ngerman sets everything for german as default hyperlinks language
\usepackage[hyphenbreaks]{breakurl} % ben�tigt f�r das Brechen von URLs in Literaturreferenzen, hyphenbreaks auch bei links, die �ber eine Seite gehen (mit hyphenation).
\usepackage{xcolor}
\definecolor{c1}{rgb}{0,0,1} % blue
\definecolor{c2}{rgb}{0,0.3,0.9} % light blue
\definecolor{c3}{rgb}{0.3,0,0.9} % red blue
\hypersetup{
    linkcolor={c1}, % internal links
    citecolor={c2}, % citations
    urlcolor={c3} % external links/urls
}
%\usepackage{cite} % needed for cite
\usepackage[square,authoryear]{natbib} % needed for cite and abbrvnat bibliography style
\usepackage[nottoc]{tocbibind} % needed for displaying bibliography and other in the table of contents
\usepackage{graphicx} % needed for \includegraphics 
\usepackage{longtable} % needed for long tables over pages
\usepackage{bigstrut} % needed for the command \bigstrut
\usepackage{enumerate} % needed for some options in enumerate
%\usepackage{todonotes} % needed for todos
\usepackage{makeidx} % needed for creating an index
\makeindex
\usepackage{gensymb}
%\usepackage{url}
\usepackage{xurl}
\usepackage{psfrag}
\usepackage{multirow}
\usepackage{subfigure}

\usepackage{algpseudocode}
\usepackage{float}
%\usepackage{minted}
\usepackage{tcolorbox}
\usepackage{menukeys}
\usepackage[
width=.8\textwidth,
justification=centering]{caption}
\usepackage{xcolor}
\usepackage{pstricks}

\usepackage[ framed, numbered]{matlab-prettifier}%framed,%
\usepackage{listings}
\usepackage{physics}
\usepackage{pdfpages}
\usepackage[toc,page]{appendix}
\usepackage{times}
\usepackage{acronym}


\usepackage{array}
%% page settings

\usepackage[top=25mm, bottom=25mm,left=25mm,right=25mm]{geometry} % needed for page border settings
\parindent=0mm % for space of first line of new text block
\sloppy % for writing with hyphenless justification (tries to)
\hyphenation{} % use hyphenation of tolerance parametershttp://www.jr-x.de/publikationen/latex/tipps/zeilenumbruch.html
\hyphenpenalty=10000
\exhyphenpenalty=10000
\usepackage{fancyhdr} % needed for head and foot options

\usepackage{setspace}
\onehalfspacing
%% my macros

%% Text fomats
\newcommand{\tbi}[1]{\textbf{\textit{#1}}}
\renewcommand\bibname{References}

% IEEE naming convention
%\renewcommand{\figurename}{Fig.}

% custom counter for shell scripts
\newcounter{shellcounter}
% \newcounter{myexample}% preamble
\newtcolorbox[
use counter=shellcounter,
number format=\Alph]{shell}[2][]{%
	colback=blue!5!white,
	colframe=blue!75!white,
	fonttitle=\bfseries,
	title= {\sc Shell \thetcbcounter}: #2,#1}


\newcounter{outputcounter}
\newtcolorbox[
use counter=outputcounter,
number format=\Roman]{result}[2][]{%
	colback=green!5!white,
	colframe=green!75!black,
	fonttitle=\bfseries,
	title= {\sc Result box \thetcbcounter}: #2,#1}
\renewcommand\bibname{References}
\usepackage{acronym}



\begin{document}
\begin{titlepage}
\center % Center everything on the page

%-------------------------------------------------------------------------------------
%	HEADING SECTIONS
%------------------------------------------------------------------------------------

\textbf{\large UNIVERSITY OF MORATUWA }\\[8mm]
{\large Faculty of Engineering}\\[1.5cm]
\includegraphics[width=0.3\textwidth]{figures/uomlogo}\\[1.5cm]


{\large Registered Module No: EN3992}\\[5mm]
\textbf{\large INDUSTRIAL TRAINING REPORT}\\[1cm]


{\large LE Robotics (Pvt.) Ltd.}\\[0.5cm]

{\large From: 04/ 01 / 2022 To: 20 / 06 / 2022}\\[1cm]

{\large Date of Submission:}\\[2mm]
{\large \today} \\[1cm]

{\large Thalagala B.P.}\\[4mm]
{\large 180631J}\\[1cm]

{\large Department of: Electronic and Telecommunication Engineering}

%----------------------------------------------------------------------------------------

\vfill % Fill the rest of the page with whitespace

\end{titlepage}
%---------------------------------------------------------------------------


\chapter*{Preface}
\addcontentsline{toc}{chapter}{Preface}
% A brief account of the Preface
This report was composed as partial fulfillment of the requirements of Module EN3992 - Industrial Training, in the curriculum of B.Sc. in Engineering (Electronic and Telecommunication) at the University of Moratuwa, Sri Lanka. The experience and knowledge that I gained during the six months of my industrial training were used and were the inspiration to create this report.\\

 

\cleardoublepage
%---------------------------------------------------------------------------



\chapter*{Acknowledgment}
\addcontentsline{toc}{chapter}{Acknowledgment}
%Appreciation of those who helped in the internship process

I would like to gratefully acknowledge all of the people who helped me to make this six months of special industrial training period a massive success, starting from the point of applying to a company to the point I left the training organization at the completion of six months.\\

First of all, I would like to express my heartfelt gratitude to Professor Kapila Jayasinghe who was our supervisor for us throughout the six months of the internship period. The advice he provided us with regard to professional engineering practices and ethics was invaluable. In addition to that the directions he provided us to gain the required technical skills required for the allocated project were priceless. I believe the mindset that he built in us, towards working under minimum supervision, will be a massive support for us to thrive in the fast-moving industry.\\

Next, I would like to express my gratitude toward Miss. Laknie Jayasinghe who was the Engineer in charge of us in our training period. The support she provided us to improve our soft skills as well as technical skills as a professional engineers, is highly appreciated. The support she provided when composing the technical documentation of my allocated task by pointing out the mistakes and the areas of improvement was invaluable and must be mentioned at this point. Moreover, her support in validating the project deliveries at the end of the training period is highly appreciated.\\

Last but not least, I would like to express my heartiest gratitude to Mr Janka Kulathunga, the manager of the Lanka Electronics research and development section and Mr Aloka Perera, an electronic engineer in LE Robotics (Pvt.) Ltd for being available for us whenever we needed their support.



%---------------------------------------------------------------------------

\tableofcontents %Three header levels (e.g. 2.7.1) are adequate
\vfill
\begin{center}
	\textbf{\textit{*PDF is clickable}}
\end{center}

%---------------------------------------------------------------------------
%Descriptions of Abbreviations used
\chapter*{List of Abbreviations}
\addcontentsline{toc}{chapter}{List of Abbreviations}
\begin{acronym}
	\acro{can}[CAN]{Controller Area Network}
	\acro{cnc}[CNC]{Computer Numerical Control}
	\acro{cpld}[CPLD]{Complex Programmable Logic Device}
	\acro{cv}[CV]{Computer Vision}
	
	\acro{fov}[FOV]{Field of View}
	\acro{fpga}[FPGA]{Field Programmable Gate Arrays}
	
	\acro{gui}[GUI]{Graphical User Interface}
	
	\acro{ide}[IDE]{Integrated Development Environment}
	\acro{ipc}[IPC]{Inter-process communication}
	
	\acro{lan}[LAN]{Local Area Network}
	
	\acro{ml}[ML]{Machine Learning}
	
%	\acro{oop}[OOP ]{Object Oriented Programming }
	\acro{oss}[OSS]{Open-Source Software}
	
	\acro{pc}[PC]{Personal Computer}
	\acro{pcb}[PCB]{Printed Circuits Boards}
	
	\acro{rnd}[R\&D]{Research and Development}
	\acro{roi}[ROI]{Region of Interest}
	\acro{rpi}[RPi]{Raspberry Pi}
			
	\acro{scara}[SCARA]{Selective Compliance Articulated Robot Arm}
	\acro{sdlc}[SDLC]{Software Development Life Cycle}
	\acro{se}[SE]{Software Engineering}
	\acro{sftp}[SFTP]{SSH(Secure shell) File Transfer Protocol}
	\acro{sift}[SIFT]{Scale Invariant Feature Transform}
	\acro{svm}[SVM]{Support Vector Machine}
	
	\acro{uart}[UART]{Universal Asynchronous Receiver/Transmitter}
\end{acronym}



%---------------------------------------------------------------------------
\listoffigures %Figures (including charts) Numbered as per the Chapter
%---------------------------------------------------------------------------
\listoftables %Tables Numbered as per the Chapter
%---------------------------------------------------------------------------




\pagebreak

\chapter{Description of the organization and business, its past, present and	future}

For my special industrial training, I got the opportunity to work as an engineering trainee at LE Robotics (Pvt.) Ltd. located at 100/4, Divulapitiya Road, Minuwangoda, Sri Lanka.It is a local Research and Development (R\&D) facility where they work on industrial articulated robot arms and related technologies. This chapter provides an extensive description of the organization and business as well as information about its past, present and future.

\section{Organization and Business}

LE robotics (Pvt.) Ltd. is a local Research and Development facility that has been established by Prof. J.A.K.S. Jayasinghe who is a senior professor in the Department of Electronic and Telecommunication Engineering at the University of Moratuwa in Sri Lanka.\\

LE robotics (Pvt.) Ltd. is the first in the Sri Lankan market to offer fully customisable robotics solutions `Made in Sri Lanka' for various automation needs for an affordable price with expertise based in Sri Lanka. In the year 2005, the company designed and manufactured the first custom robotics solution in their affiliated company, Lanka Electronics (Pvt.) Ltd. Since then, they have been developing various robotics solutions and related technologies in the facility.

\begin{figure}[H]
	\centering
	\fbox{
		\includegraphics[width=0.5\linewidth]{figures/logoler}
	}
	\caption{Logo of the LE Robotics (Pvt.) Ltd.}
	\label{fig:logoler}
\end{figure}

Following are a few of the features that the LE Robotics (Pvt.) Ltd. has to offer its customers.

\begin{itemize}
	\item Ease of robot customization to suit your process needs
	\item Various modes of operation depending on various motion and speed requirements
	\item Simple user interface with multi device support
	\item Convenient availability of spare parts
	\item Electronic control systems tailor-made to your organizational requirements
	\item Superior quality and reliability
\end{itemize}


\chapter{Description of familiarization work carried out}

\chapter{Exposure to systems (HSE, Financial, Administration, Logistics, etc.)}

%===========================================================================
\chapter{Project Work}

By title the project that I was assigned, was \textbf{\textit{Machine vision based Real time Motion Planning for an Industrial Articulated Robot Arm}}. In simple words it was a project related to an automatic pick and place machine which can be used to pick objects placed on a conveyor belt and pass them to the next stage of processing.\\

My contribution to that project was to develop the following three aspects of the system. 

\begin{enumerate}[I.]
	\item Development of an object detection framework: 
	\item Development of an application to train an object classification model
	\item Development of an application for camera calibration
\end{enumerate}

Subsequent sections will thoroughly explain the mentioned sub projects that were undertaken by me.

\section{Development of an Object Detection Framework}

The first project that I was assigned as a trainee electronic engineer was related to Computer vision field. An Object Detection Framework was developed to be deployed in an Automatic Pick and Place Machine. The framework is capable of, identifying regions of interests (ROIs), detecting and classifying objects, determining location and orientation of objects with respect to a real world coordinate system for grasping (picking).

%===========================================================================
\chapter{Hands on experiences}

\chapter{Soft Skills Development}

\chapter{SWOT Analysis of the organization and self}

\chapter{Conclusion: Own perspective of areas to be improved (of the whole training process including self)}


\begin{appendices}
	\chapter{Guidelines}
\end{appendices}

\bibliographystyle{plain}
\bibliography{refer}

%---------------------------------------------------------------------------
\end{document}
-
